% GENERATED FILE. Edit canonical lessons, then run npm run generate:books.
% canonical-source-hash: fnv1a64:822a60cd3861a950
% canonical-lessons: LA-C285-qua-via, LA-C285-recta-via, LA-C285-iter, LA-R285-asking-the-way

\chapter{Asking the Way}
\label{ch:la-asking-the-way}
\hlchaptermodality{285}

\begin{chapteropening}
\textbf{By the end of this chapter:} \emph{I can ask the way politely, follow directions to the right, to the left and straight on, and talk about a journey.}

Complete LA-R285-asking-the-way: I can ask the way politely, follow directions to the right, to the left and straight on, and talk about a journey.
\end{chapteropening}

\section[quā viā ad forum?]{quā viā — asking the way}

\label{lesson:LA-C285-qua-via}



\begin{quote}
Ask \textquotedblleft{}where is the forum?\textquotedblright{} (\emph{Ubi est forum?})

That asks for a place. A traveller standing at a crossroads needs something else: not where it is, but which way to go.
\end{quote}

\subsection*{You'll want to know: via, viae}
\textbf{via, viae} — \textquotedblleft{}a road, a way, a street\textquotedblright{}. It is a feminine noun like \emph{porta}.

You own \textbf{quā}, \textquotedblleft{}by which way\textquotedblright{}, and its own lesson said that \emph{viā}, \textquotedblleft{}road\textquotedblright{}, was understood inside it. Now say it out loud:

\begin{quote}
\textbf{Quā viā ad forum?} — By which road to the forum?
\end{quote}

\emph{Viā} is the ablative, \textquotedblleft{}by road\textquotedblright{}: the road you travel along.

\begin{grammarlens}[title={Grammar Lens: asking politely}]
A stranger deserves a softer opening. Put two pieces in front:

\begin{quote}
\textbf{Quaesō, dīc mihi: quā viā ad thermās?} — Please, tell me: which way to the baths.
\end{quote}

\textbf{Dīc} is \textquotedblleft{}tell!\textquotedblright{}, the command of \emph{dīcō}, and it is short: no ending at all. \emph{Mihi}, \textquotedblleft{}to me\textquotedblright{}, you know from \emph{ignōsce mihi}.

\emph{Ad} plus the accusative is where you are heading: \emph{ad forum}, \emph{ad thermās}.
\end{grammarlens}

\begin{cousinweb}
\textbf{via} is inherited, usually traced to Proto-Indo-European \emph{*wegh-} (EtymDB 2.1), the root the \emph{vehō} lesson named, \textquotedblleft{}to carry, to convey\textquotedblright{}; Lewis \& Short put the two words together. A road is where things are carried. English \emph{way} is usually given as a cousin, and English \emph{via}, \textquotedblleft{}by way of\textquotedblright{}, is this word's ablative, taken over whole.
\end{cousinweb}

\subsection*{Your turn}
Say these aloud:

\begin{itemize}
\raggedright
  \item \emph{via, viae}
  \item \emph{Quā viā ad forum?}
  \item \emph{Quaesō, dīc mihi: quā viā ad thermās?}
\end{itemize}

\begin{tcolorbox}[breakable,colback=teal!4,colframe=teal!35!black,title={Before you move on}]
Say \textquotedblleft{}a road\textquotedblright{}, then \textquotedblleft{}by which road to the forum\textquotedblright{}. (\emph{\textbf{Via. Quā viā ad forum.}}) Say \textquotedblleft{}tell me!\textquotedblright{} (\emph{\textbf{Dīc mihi!}}) Which verb you own is \emph{via} related to? (\emph{\textbf{Vehō}}, to carry by cart.)
\end{tcolorbox}
\section[rēctā viā, ad dextram, ad sinistram]{rēctā viā, ad dextram — following the directions}

\label{lesson:LA-C285-recta-via}



\begin{quote}
Ask the way to the forum, politely. (\emph{Quaesō, dīc mihi: quā viā ad forum?}) And say \textquotedblleft{}right\textquotedblright{}, then \textquotedblleft{}left\textquotedblright{}. (\emph{Dexter. Sinister.})
\end{quote}

\begin{grammarlens}[title={Grammar Lens: three directions}]
\begin{quote}
\textbf{Rēctā viā ī.} — Go straight on. \textbf{Deinde ad dextram.} — Then to the right. \textbf{Tum ad sinistram.} — Then to the left.
\end{quote}

\textbf{ī} is the command of \emph{eō}: \textquotedblleft{}go!\textquotedblright{}, a whole sentence in one long vowel. \emph{Rēctā viā} is \textquotedblleft{}by the straight road\textquotedblright{}, the ablative of \emph{via} with \emph{rēctus} agreeing.

\textbf{ad dextram} and \textbf{ad sinistram} are \textquotedblleft{}towards the right\textquotedblright{} and \textquotedblleft{}towards the left\textquotedblright{}: \emph{dextram} and \emph{sinistram} are the feminine accusative, because \emph{manum}, \textquotedblleft{}hand\textquotedblright{}, is understood. Where something already is, Latin says \textbf{dextrā}, \textquotedblleft{}on the right\textquotedblright{}, and \textbf{sinistrā}, \textquotedblleft{}on the left\textquotedblright{}.

\begin{quote}
\textbf{Thermae sunt dextrā, iūxtā pontem.} — The baths are on the right, next to the bridge.
\end{quote}
\end{grammarlens}

\begin{culture}
Roman comedy has a scene built on exactly this. In Terence's \emph{Adelphoe}, a slave gets rid of an old man by sending him across town: \emph{ad sinistram}, then \emph{ad dextram}, along \emph{hāc rēctā plateā}, \textquotedblleft{}this straight street\textquotedblright{}, to a place where the person he wants is not. Lewis \& Short cite \emph{ad dexteram}, the longer spelling, from Plautus and from Terence's \emph{Andria}. A Roman audience laughed because they gave and got directions this way every day.
\end{culture}

\begin{cousinweb}
\textbf{dextram} is the adjective \emph{dexter} with \emph{manum} left out: Lewis \& Short take \emph{dextra}, \textquotedblleft{}the right\textquotedblright{}, as \emph{dextra manus}, \textquotedblleft{}the right hand\textquotedblright{}, with the noun understood. The \emph{dexter} lesson traced the word itself to Proto-Indo-European \emph{*deks-}, \textquotedblleft{}right\textquotedblright{}, and named Greek \emph{dexios} as a cousin.
\end{cousinweb}

\subsection*{Your turn}
Say these aloud:

\begin{itemize}
\raggedright
  \item \emph{Rēctā viā ī.}
  \item \emph{Prīmum rēctā viā ī; deinde ad dextram.}
  \item \emph{Tum ad sinistram. Ibi est forum.}
  \item \emph{Thermae sunt dextrā, iūxtā pontem.}
  \item directions of your own, from your door to the nearest shop
\end{itemize}

\begin{tcolorbox}[breakable,colback=teal!4,colframe=teal!35!black,title={Before you move on}]
\textquotedblleft{}Go straight on\textquotedblright{}? (\emph{\textbf{Rēctā viā ī.}}) \textquotedblleft{}To the left\textquotedblright{}? (\emph{\textbf{Ad sinistram.}}) \textquotedblleft{}On the right\textquotedblright{}? (\emph{\textbf{Dextrā.}}) What word is understood inside \emph{ad dextram}? (\emph{\textbf{Manum}}, hand.)
\end{tcolorbox}
\section[iter, itineris]{iter, itineris — a journey}

\label{lesson:LA-C285-iter}



\begin{quote}
Say \textquotedblleft{}go straight on\textquotedblright{}. (\emph{Rēctā viā ī.}) Then \textquotedblleft{}we go\textquotedblright{}. (\emph{Īmus.})
\end{quote}

\subsection*{You'll want to know: iter, itineris}
\textbf{iter, itineris} — \textquotedblleft{}a journey; a route, a way\textquotedblright{}. It is neuter, and the genitive shows the stem every other form is built on: \textbf{itiner-}.

\begin{quote}
\textbf{Iter longum erat.} — The journey was long. \textbf{In itinere} — on the way, on the journey
\end{quote}

\emph{Via} is the road under your feet; \emph{iter} is the going along it, from the start to the arrival.

\subsection*{You'll want to know: iter faciō}
To travel, a Roman \textbf{made a journey}: \textbf{iter faciō}. Caesar writes it of whole armies marching through a province.

\begin{quote}
\textbf{Rōmam iter faciō.} — I am travelling to Rome. \textbf{Herī iter fēcī.} — Yesterday I made a journey.
\end{quote}

\emph{Rōmam} needs no \emph{ad}: a city's name points the way by itself, as \emph{domum} does. \emph{Fēcī} is the perfect of \emph{faciō}, a stem learned on its own.

\begin{cousinweb}
\textbf{iter} is built on \emph{īre}, \textquotedblleft{}to go\textquotedblright{} (Lewis \& Short), the one-vowel verb whose lesson traced it to Proto-Indo-European \emph{*h\textsubscript{1}ey-}: a journey is a going. Its stem \emph{itiner-} gives English \emph{itinerary}, a plan of the route, and \emph{itinerant}, someone who lives on the road.
\end{cousinweb}

\subsection*{Your turn}
Say these aloud:

\begin{itemize}
\raggedright
  \item \emph{iter, itineris}
  \item \emph{Iter longum erat.}
  \item \emph{Rōmam iter faciō.}
  \item \emph{Herī iter fēcī.}
  \item the road and the journey — \emph{via}, then \emph{iter}
\end{itemize}

\begin{tcolorbox}[breakable,colback=teal!4,colframe=teal!35!black,title={Before you move on}]
What does \emph{iter, itineris} mean? (\textbf{A journey; a route.}) How does a Roman say \textquotedblleft{}I travel\textquotedblright{}? (\emph{\textbf{Iter faciō.}}) And which verb is inside \emph{iter}? (\emph{\textbf{Īre}}, to go.)
\end{tcolorbox}
\section[(review)]{Review — at the crossroads}

\label{lesson:LA-R285-asking-the-way}



\begin{quote}
Say \textquotedblleft{}a road\textquotedblright{}, then \textquotedblleft{}a journey\textquotedblright{}. (\emph{Via. Iter.})
\end{quote}

\begin{grammarlens}[title={What you've built}]
Three lessons, and you can find your way through a Roman town. The question: \emph{quaesō, dīc mihi: quā viā ad …?} The answer: \emph{rēctā viā ī}, \emph{ad dextram}, \emph{ad sinistram}, and \emph{dextrā} or \emph{sinistrā} for where a place already stands. And the whole of it: \emph{iter faciō}, I am on a journey.

\begin{quote}
\emph{— Quaesō, dīc mihi: quā viā ad thermās?} \emph{— Prīmum rēctā viā ī. Deinde ad sinistram. Thermae sunt dextrā, iūxtā pontem.}
\end{quote}

The rungs of a story, \emph{prīmum} and \emph{deinde}, are the rungs of a set of directions too.
\end{grammarlens}

\subsection*{Your turn}
Say these aloud:

\begin{itemize}
\raggedright
  \item the question — \emph{Quaesō, dīc mihi: quā viā ad thermās?}
  \item the answer — \emph{Prīmum rēctā viā ī. Deinde ad sinistram.}
  \item \emph{Thermae sunt dextrā, iūxtā pontem.}
  \item advice for the road — \emph{Rēctā viā īre dēbēs, quia via longa est.}
  \item the advice too late — \emph{Viam rogāre dēbēbās.} — You should have asked the way
  \item ask the way to the forum, then give the answer yourself
\end{itemize}

\begin{tcolorbox}[breakable,colback=teal!4,colframe=teal!35!black,title={Before you move on}]
\textquotedblleft{}Which way to the baths?\textquotedblright{} (\emph{\textbf{Quā viā ad thermās?}}) \textquotedblleft{}Go straight on\textquotedblright{}? (\emph{\textbf{Rēctā viā ī.}}) \textquotedblleft{}To the right\textquotedblright{}? (\emph{\textbf{Ad dextram.}}) \textquotedblleft{}I am travelling to Rome\textquotedblright{}? (\emph{\textbf{Rōmam iter faciō.}})
\end{tcolorbox}
